\documentclass[11pt,a4paper]{jlreq}
\usepackage{amsmath,amssymb}
\usepackage{graphicx}

\begin{document}

\begin{center}
{\LARGE \textbf{情報科学基礎A 第6回 レポート}}
\end{center}
\vspace{1em}

\begin{flushright}
提出日：\today \\
学籍番号：2611193 \\
氏名：中井平蔵
\end{flushright}

\section*{Problem 1}
\begin{enumerate}
    \item[(1)]\leavevmode\\
    1枚のコインを1回投げる試行において、結果を「表裏」だけでなく、自分の体を基準とした左右位置でも分類する。ここで $H$ は表、$T$ は裏を表す。
    立ち位置を変えずに、体の真ん中を通る線を基準に、左側($L$)、右側($R$)、およびその線の上($M$)の3つに分ける。
    したがって、標本空間 $\Omega$ を以下のように定義する。
    \[
    \fbox{$\Omega = \{(H,L), (H,R), (H,M), (T,L), (T,R), (T,M)\}$}
    \]

    \item[(2)]\leavevmode\\
    次に、各標本に対して確率を定義する。
    初めに、体の真ん中を通る線の上に落ちる確率を
    \[
    P(M)=p
    \]
    とすると、線の上ではない（左または右）確率は、
    \[
    P(L\cup R)=1-p
    \]
    となる。左右対称より、
    \[
    P(L)=P(R)=\frac{1-p}{2}
    \]
    となる。ここで、コインの表裏はコインの位置と独立で、かつ
    \[
    P(H)=P(T)=\frac12
    \]
    より、
    \[
    P(H,L)=P(H)P(L)=\frac12\cdot\frac{1-p}{2}=\frac{1-p}{4}
    \]
    同様に、
    \[
    P(H,R)=P(T,L)=P(T,R)=\frac{1-p}{4}
    \]
    また、
    \[
    P(H,M)=P(T,M)=\frac12\cdot p=\frac{p}{2}
    \]
    したがって、各標本に対して確率を以下のように定義する。
    \[
    \fbox{$P(H,L)=P(H,R)=P(T,L)=P(T,R)=\frac{1-p}{4},\quad P(H,M)=P(T,M)=\frac{p}{2}$}
    \]
\end{enumerate}

\section*{Problem 2}

\begin{enumerate}
\item[(1)]\leavevmode\\

与えられた分布関数より、
\[
\fbox{$P(X \le 0.5) = 0.5$}
\]

\item[(2)]\leavevmode\\

確率密度関数を $f(x)$ とすると、確率は以下のように表せる。

\[
P(a\le X\le b)=\int_a^b f(x)\,dx
\]

ここで、一点 $X=0.5$ における確率$P(X=0.5)$は、以下のように考えることができる。

\[
P(X=0.5)
=
P(0.5 \leq X \leq 0.5)
\]

よって、

\[
P(X=0.5)
=
\int_{0.5}^{0.5} f(x)\,dx
\]

となる。

積分区間の幅が $0$ であるため、

\[
\int_{0.5}^{0.5} f(x)\,dx = 0
\]

である。したがって、

\[
\fbox{$P(X=0.5)=0$}
\]
\end{enumerate}


\section*{Problem 3}
\noindent
宝くじの平均獲得金額を求めることは、1回あたりの獲得金額の期待値 $E[X]$ を求めることである。
\noindent
宝くじの総発行枚数を $N$、$i$ 等の当選本数を $n_i$ とすると、各賞の当選確率は
\[
P(X=x_i)=\frac{n_i}{N}
\]

\noindent で与えられる。
また、販売枚数 $M$ が総発行枚数 $N$ 以下の場合でも、券が無作為に売れると考えれば、購入した1枚が $i$ 等に当たる確率は同じく $n_i/N$ とみなせる。
したがって、各賞の当選金額を $x_i$ とすると、期待値は以下のようになり、平均獲得金額を推定できる。
\[
E[X] = \sum_i x_i P(X=x_i)
\]

\noindent
しかし、分散が非常に大きい宝くじでは、上記の期待値に収束するために十分大きな試行回数が必要である。
これは、分散が大きいほど1回ごとの獲得金額の変動が大きく、少ない試行回数では標本平均が期待値から大きく乖離しやすいためである。
そのため、期待値に収束するためには、より多くの試行回数を重ねる必要がある。したがって、

\[
\fbox{\parbox{0.9\linewidth}{
期待値 $E[X]=\sum_i x_iP(X=x_i)$ を計算することで平均獲得金額を推定できる。
しかし、分散が非常に大きい宝くじの場合、少ない試行回数では期待値に収束せず、平均獲得金額を推定することは難しい。
十分に多くの試行回数を重ねることで期待値に収束する。
}}
\]

\end{document}
